\chapter{Definición y Planteamiento del Problema}

\section{Planteamiento del Problema}

El entorno productivo de los talleres industriales enfrenta diversas problemáticas que afectan su eficiencia y rendimiento. Entre los principales desafíos se encuentran la capacidad de producción limitada, los largos tiempos de fabricación y los retrasos en las fechas de entrega \cite{ref7}. Además, la falta de sistemas estructurados para la evaluación de calidad y la baja integración tecnológica pueden impactar negativamente la productividad del proceso \cite{ref7}.  

La asignación de tareas en los talleres industriales es un factor clave para optimizar el uso de recursos humanos y técnicos. Una planificación ineficiente puede provocar retrasos, una distribución desequilibrada de la carga de trabajo y una reducción en la productividad \cite{ref8}. La complejidad de estos entornos productivos, caracterizados por la variabilidad en la demanda y la necesidad de coordinar múltiples operaciones, hace necesario el uso de enfoques avanzados que permitan gestionar eficazmente la distribución de actividades y mejorar el cumplimiento de los plazos de entrega \cite{ref8}.  

Para abordar este problema, se propone la implementación de algoritmos genéticos (AG) en un módulo ERP de gestión de órdenes de trabajo. Esta solución permitirá automatizar la distribución de tareas considerando factores clave como:

\begin{itemize}
    \item Capacidades técnicas de cada operario.
    \item Disponibilidad de insumos requeridos para cada tarea.
    \item Carga de trabajo actual y disponibilidad horaria de los operarios.
\end{itemize}

Adicionalmente, este estudio evaluará el desempeño del algoritmo genético en comparación con otro método de optimización utilizado en ingeniería industrial, con el objetivo de determinar su viabilidad y eficacia en este contexto.


\subsection{Formulación}
¿Cómo mejorar la asignación de tareas en un ERP de talleres industriales utilizando algoritmos genéticos, y cómo se compara su rendimiento con otras técnicas de automatización utilizadas en ingeniería industrial?

\subsection{Sistematización}
Para abordar de manera integral el problema de investigación, se plantean las siguientes preguntas:

\begin{enumerate}
    \item ¿Cuáles son los principales factores que afectan la asignación de tareas en talleres industriales?
    \item ¿Cómo pueden los algoritmos genéticos mejorar la automatización de tareas dentro del ERP?
    \item ¿Qué otras técnicas de automatización existen en ingeniería industrial para resolver problemas similares?
    \item ¿Cómo se comparan los resultados obtenidos con algoritmos genéticos respecto a otra técnica de automatización en términos de eficiencia y calidad de la solución?
\end{enumerate}

\section{Objetivos}

\subsection{Objetivo General} 
Desarrollar un modelo de automatización basado en algoritmos genéticos para apoyar en el módulo de asignación de tareas de un sistema ERP de talleres industriales, comparándolo con otro método de automatización utilizado en ingeniería industrial.

\subsection{Objetivos Específicos}
Para alcanzar el objetivo general, se plantean los siguientes objetivos específicos:

\begin{enumerate}
    \item Identificar las caracaterísticas claves que influyen en la asignación de tareas en talleres industriales.
    \item Diseñar e implementar un modelo utilizando un algoritmo genético para automatizar la asignación de tareas tales como desmontaje del motor, limpieza y diagnóstico, rectificación de cilindros, rectificación de cigüeñales, ajuste de válvulas, ensamblaje y prueba del motor.
    \item Evaluar la funcionalidad del modelo y el desempeño del algoritmo genético en comparación con otra técnicas de automatización utilizadas en ingeniería industrial.
    \item Analizar los beneficios y limitaciones de los algoritmos genéticos en este contexto.
\end{enumerate}

\section{Justificación}

La gestión eficiente de órdenes de trabajo en talleres industriales impacta directamente la productividad y rentabilidad. El uso de algoritmos genéticos puede ofrecer una solución flexible y adaptable que automatiza los tiempos de ejecución, minimiza costos y mejora la distribución equitativa de la carga de trabajo.  

La empresa Rectificadora Recti-Ram  cuenta con datos históricos sobre tiempos de reparación, asignación de operarios y disponibilidad de recursos, lo que facilita el entrenamiento y validación del modelo propuesto. Asimismo, la disponibilidad de la empresa para proporcionar información clave sobre su flujo de trabajo y sus restricciones operativas permitirá diseñar un modelo que se ajuste a las necesidades reales.  

Además, comparar esta solución con otro método de automatización permitirá validar su aplicabilidad en escenarios industriales reales, asegurando que la mejor estrategia sea implementada dentro del ERP del taller, garantizando así una planificación eficiente y adaptable a las condiciones dinámicas del entorno productivo.

\section{Delimitaciones y Alcances}

\begin{itemize}
    \item Se enfocará en la asignación de tareas dentro del módulo de ordenes de trabajo en un ERP para talleres industriales.
    \item Se implementará un modelo basado en algoritmos genéticos.
    \item Se comparará con otro método de automatización (a definir en consulta con expertos en ingeniería industrial).
\end{itemize}